\documentclass[a4paper, oneside, 12pt]{report}

\usepackage[T1]{fontenc}
\usepackage{libertinus}
\usepackage{underscore}
\usepackage{geometry}
\usepackage{float}
\usepackage{titling}
\usepackage{titlesec}
\usepackage{paralist}
\usepackage{footnote}
\usepackage{enumitem}
\setitemize{noitemsep,topsep=0pt,parsep=0pt,partopsep=0pt}
\usepackage{amsmath, amsthm}
\usepackage{gb4e}
\noautomath
\usepackage{bbm}
\usepackage{textcomp}
\usepackage{soul}
\usepackage{graphicx}
\usepackage{siunitx}
\usepackage[table,xcdraw]{xcolor}
\usepackage{tikz}
\usepackage[ruled, vlined, linesnumbered, noend]{algorithm2e}
\usepackage{xr-hyper}
\usepackage[colorlinks, citecolor = purple, bookmarksnumbered]{hyperref} % linkcolor=black, anchorcolor=black, citecolor=black, filecolor=black
\usepackage[most]{tcolorbox}
\usepackage{caption}
\usepackage{subcaption}
\usepackage{booktabs}
\usepackage{multirow}
\usepackage[figuresright]{rotating}
\usepackage{acro}
\usepackage[round]{natbib} 
\usepackage{nameref,zref-xr}
\zxrsetup{toltxlabel}
\zexternaldocument*[alignment-]{../alignment/alignment}[alignment.pdf]
\zexternaldocument*[exercise1-]{../Exercise/2021-3}[2021-3.pdf]
\zexternaldocument*[method-]{../methodology/glossing}[glossing.pdf]
\usepackage{prettyref}

\geometry{left=3.18cm,right=3.18cm,top=2.54cm,bottom=2.54cm}
\titlespacing{\paragraph}{0pt}{1pt}{10pt}[20pt]
\setlength{\droptitle}{-5em}

\DeclareMathOperator{\timeorder}{\mathcal{T}}
\DeclareMathOperator{\diag}{diag}
\DeclareMathOperator{\legpoly}{P}
\DeclareMathOperator{\primevalue}{P}
\DeclareMathOperator{\sgn}{sgn}
\newcommand*{\ii}{\mathrm{i}}
\newcommand*{\ee}{\mathrm{e}}
\newcommand*{\const}{\mathrm{const}}
\newcommand*{\suchthat}{\quad \text{s.t.} \quad}
\newcommand*{\argmin}{\arg\min}
\newcommand*{\argmax}{\arg\max}
\newcommand*{\normalorder}[1]{: #1 :}
\newcommand*{\pair}[1]{\langle #1 \rangle}
\newcommand*{\fd}[1]{\mathcal{D} #1}
\newcommand*{\textto}{$\to$}
\newcommand*{\textgt}{$>$ }

\newcommand*{\citesec}[1]{\S~{#1}}
\newcommand*{\citechap}[1]{Ch~{#1}}
\newcommand*{\citefig}[1]{Fig.~{#1}}
\newcommand*{\citetable}[1]{Table~{#1}}
\newcommand*{\citepage}[1]{p.~{#1}}
\newcommand*{\citepages}[1]{pp.~{#1}}
\newcommand*{\citefootnote}[1]{fn.~{#1}}
\newcommand*{\citechapsec}[2]{\citechap{#1}.\citesec{#2}}

\newrefformat{sec}{\citesec{\ref{#1}}}
\newrefformat{fig}{\citefig{\ref{#1}}}
\newrefformat{tbl}{\citetable{\ref{#1}}}
\newrefformat{chap}{\citechap{\ref{#1}}}
\newrefformat{fn}{\citefootnote{\ref{#1}}}
\newrefformat{box}{Box~\ref{#1}}
\newrefformat{ex}{\ref{#1}}


% color boxes

\tcbuselibrary{skins, breakable, theorems}

\AtBeginEnvironment{infobox}{\small}
\AtBeginEnvironment{theorybox}{\small}

\newtcbtheorem[number within=chapter]{infobox}{Box}{
    enhanced,
    boxrule=0pt,
    colback=blue!5,
    colframe=blue!5,
    coltitle=blue!50,
    borderline west={4pt}{0pt}{blue!65},
    sharp corners,
    fonttitle=\bfseries, 
    breakable,
    before upper={\parindent15pt\noindent}}{box}
\newtcbtheorem[number within=chapter, use counter from=infobox]{theorybox}{Box}{
    enhanced,
    boxrule=0pt,
    colback=orange!5, 
    colframe=orange!5, 
    coltitle=orange!50,
    borderline west={4pt}{0pt}{orange!65},
    sharp corners,
    fonttitle=\bfseries, 
    breakable,
    before upper={\parindent15pt\noindent}}{box}
\newtcbtheorem[number within=chapter, use counter from=infobox]{learnbox}{Box}{
    enhanced,
    boxrule=0pt,
    colback=green!5,
    colframe=green!5,
    coltitle=green!50,
    borderline west={4pt}{0pt}{green!65},
    sharp corners,
    fonttitle=\bfseries, 
    breakable,
    before upper={\parindent15pt\noindent}}{box}
\newtcbtheorem[number within=chapter, use counter from=infobox]{todobox}{Box}{
    enhanced,
    boxrule=0pt,
    colback=red!5,
    colframe=red!5,
    coltitle=red!50,
    borderline west={4pt}{0pt}{red!65},
    sharp corners,
    fonttitle=\bfseries, 
    breakable,
    before upper={\parindent15pt\noindent}}{box}

% Shorthands
\newcommand*{\concept}[1]{\textbf{#1}}
\newcommand*{\term}[1]{\emph{#1}}
\newcommand{\form}[1]{\emph{#1}}

\newcommand{\redp}{\textasciitilde}

\newcommand{\deictictime}{T$_{\text{d}}$}
\newcommand{\referredtime}{T$_{\text{r}}$}
\newcommand{\orientationtime}{T$_{\text{o}}$}

\DeclareAcronym{blt}{short = BLT, long = Basic Linguistic Theory}
\DeclareAcronym{cgel}{short = CGEL, long = The Cambridge Grammar of the English Language}
\DeclareAcronym{dm}{short = DM, long = Distributed Morphology}
\DeclareAcronym{tag}{long = Tree-adjoining grammar, short = TAG}
\DeclareAcronym{sfp}{long = sentence-final particle, short = sfp}
\DeclareAcronym{np}{long = noun phrase, short = NP}
\DeclareAcronym{vp}{long = verb phrase, short = VP}
\DeclareAcronym{pp}{long = preposition phrase, short = PP}
\DeclareAcronym{advp}{long = adverb phrase, short = AdvP}
\DeclareAcronym{cls}{long = classifier, short = CLS}
\DeclareAcronym{dist}{long = distal, short = DIST}
\DeclareAcronym{prox}{long = proximate, short = PROX}
\DeclareAcronym{dem}{long = demonstrative, short = DEM}
\DeclareAcronym{classify}{long = classifier, short = \textsc{cl}}
\DeclareAcronym{dur}{long = durative, short = DUR}
\DeclareAcronym{neg}{long = negative, short = \textsc{neg}}
\DeclareAcronym{cc}{long = copular complement, short = CC}
\DeclareAcronym{cs}{long = copular subject, short = CS}
\DeclareAcronym{tame}{long = {tense, aspect, mood, evidentiality}, short = TAME}
\DeclareAcronym{past}{long = past, short = PST}
\DeclareAcronym{nonpast}{long = non-past, short = NPST}
\DeclareAcronym{present}{long = present, short = PRES}
\DeclareAcronym{progressive}{long = progressive, short = \textsc{poss}}
\DeclareAcronym{perfect}{long = perfect, short = \textsc{perf}}
\DeclareAcronym{passive}{long = passive, short = \textsc{pass}}
\DeclareAcronym{copula}{long = copula, short = COP}
\DeclareAcronym{possessive}{long = possessive, short = \textsc{poss}}
\DeclareAcronym{coca}{long = Corpus of Contemporary American English, short = COCA}
\DeclareAcronym{sap}{long = speech act participant, short = SAP}

\newcommand{\asis}[1]{\textsc{#1}}
\newcommand{\oneof}[1]{{#1}}
\newcommand*{\homo}[2]{#1$_{\text{#2}}$}
\newcommand{\category}[1]{\textsc{#1}}
\newcommand{\formcat}[1]{\textsc{#1}}
\newcommand{\emptymorpheme}{$\emptyset$}
\newcommand*{\fromto}[2]{\langle {#1}, {#2} \rangle}

\newcommand{\alignment}{\href{../alignment/alignment.pdf}{my notes about alignment}}
\newcommand{\exerciseone}{\href{../Exercise/2021-3.pdf}{this exercise}}
\newcommand{\method}{\href{../methodology/glossing.pdf}{this note about my understanding of descriptive grammars}}

\newcommand{\ala}{à la}
\newcommand{\translate}[1]{`#1'}
\newcommand{\vP}{\textit{v}P}

% Make subsubsection labeled
\setcounter{secnumdepth}{4}
\setcounter{tocdepth}{4}
% reset example counter every chapter (but do not include the chapter number to the label)
\counterwithin{exx}{chapter} 

% Reference formats
\renewcommand{\bibname}{References}
\setcitestyle{aysep={}} 

% List format
\setlist[enumerate,1]{label=\alph*\upshape)}

\title{Arabic grammar}
\author{Jinyuan Wu}

\begin{document}

\maketitle

\chapter{Introduction}

\section{History and classification}

\subsection{Classification}

Arabic -- as a collective term for all Arabic languages or ``dialects'' --
is a Central Semitic group.

It is sometimes claimed that Quranic Arabic
(which itself is not a trivial concept; \prettyref{sec:introduction.quran})
is the shared ancestor of modern dialects.
It is indeed true that 
given the fact that most Arabic dialects originated from the spoken language 
of early Islam conquerors, Old Hijazi was likely the shared ancestor 
of \emph{most} materials we see in modern dialects.
However, contrary to what many Western scholars have assumed,
linguistic diversity within the Arabic language(s)
long before the emergence of modern dialects,
a fact that had already been noticed by native Arabic grammarians
\citep[\citepages{15-39}]{van2022quranic}.
This indicates that it is probably not appropriate to assume that all strata of modern dialects are from ``Quranic Arabic''.%
\footnote{
Further, certain modern dialects near the border between Yemen and Saudi Arabia
retain certain archaic features which are not even seen in the Quran,
the most prominent one being arguably nasal definite articles,
although whether this is due to language mixing with Modern South Arabian languages is not clear at the moment \citep{watson2018south}.
}
Indeed, we will also see that ``Quranic Arabic'' itself came from admixture.

A Neogrammarian analysis of the genetic relations 
between modern Arabic dialects is current lacking.
There exist a tentative division between Western and Eastern dialects,
which is determined by verb conjunction,
roughly corresponding to the geographical border between 
the Western and Eastern part of Roman Empire.

\subsection{Quran and its language}\label{sec:introduction.quran}

The Quran is often praised to be the standard of Classical Arabic,
but there are some subtleties regarding what exactly the term \term{Quranic Arabic} means.

It is known by native Arab grammarians that the Quran shows distinct features
setting it apart from classical Arabic poetry
\citep[\citepage{4}]{van2022quranic}.
The description inside of the Quran implies that its language was the local vernacular 
(hence ``intelligible''), which was Old Hijazi \citep[\citepages{216-217}]{van2022quranic},
and the Quran consonantal skeleton, upon comparison with Arab grammarians
\citep[\citepages{117-118,145}]{van2022quranic}.

On the other hand, the reading traditions of Quran are not transparent reflections of 
any of the Arabic dialects at Muhammad's time.

First, it is likely that the reading traditions originated as 
ways of recitation of a more or less fixed written text,
and \emph{not} oral transmission.
All canonical readings agree with the Uthmanic consonantal skeleton text,
which is either because the oral traditions have a shared origin
which happened to be put down by the Uthmanic edition,
or because the oral traditions came from parsing a written consonantal skeleton text.
Because there exist discrepancies between oral traditions with rather different phonetics
but identical consonantal skeletons,
and the stemma of the readers is consistent with the stemma of Quranic manuscripts
\citep[\citepages{53-55}]{van2022quranic}.

Second, no regular sound changes can be established by comparative studies,
which is not the case for early Arabic dialects according to the Arab grammatical tradition
\citep[\citepage{55}]{van2022quranic}.
Mixture of dialectal features is common \citep[e.g. \citepage{71}]{van2022quranic}; 
some differences between the reading traditions are feasibly explained
by synchronic morphological rules but not diachronic rules (the latter having no access to the morphological structure of words),
which however were never completely attested in dialects
\citep[e.g. \citepages{83-84}]{van2022quranic}.

Therefore, the reading traditions should not be taken as transparent reflections of natural dialects 
\citep[\citepages{97-98}]{van2022quranic}.

The most popular Quran reading tradition, Hafs 'an 'Asim, is very regular. 
It still likely does not represent any real historic dialect.
Many features of it were probably due to hypercorrection,
like elimination of certain vowels 
(likely due to archaism based on a real dialect of a nomad tribe
which indeed contained only three vowels).
Certain quirky features sporadically appear,
but they can also be conceived as the result of early scholars showing off their knowledge 
of (perceived) ancient features of Arabic dialects of their times.
However, this of course is not contradictory to the assumption that 
certain reading conventions do possibly reflect some features of the spoken language 
at the time of Muhammad.



\subsection{The makeup of Classical Arabic}


What is known as Classical Arabic is a mixture of Old Hijazi 

\chapter{Grammatical overview}

\paragraph*{The construct state}
It seems the construct state construction is not too far from the English \form{a [field linguist]}:
the modifier in the construction is not a full, ``sealed'' \ac{np},
just like the word \form{field} is a \term{nominal} and not a full \ac{np}.
The main differences are 
(a) in Arabic the head noun always appears initially and 
(b) the article is shared by all nominals \citep{alqahtani2016structure}, and 
(c) most importantly, the ``modifier'' in a construct phrase can be referential.

1. Can the head in the genitive construction receive the definite marker? (13) seems strange?
2. Definite head with indefinite possessor? (21) seems wrong: the wrong sequence \emph{can} be understood as a definite NP
3. Indefinite head with definite possessor and indefinite attributives (57b) seems wrong

A possible reason for ``definite spreading'': head-(INDEF, dropped) possessor-DEF adj-INDEF is often regarded as a nominal predication construction.
Why this parsing is preferred however remains to be answered.

Definiteness spreading per se seem comparable to the English \form{the country's king's properties},
which is rendered definite as a whole because of the subject-possessor \form{the country}.
The problem is the definiteness of this NP can probably explained by semantic reasons,
whereas in Semitic languages, it seems \form{king} and \form{properties} each has a \emph{syntactic} definiteness value,
as is shown by the definiteness agreement between these words and adjectives modifying them.
This is not a very local agreement because the construct state chain can be arbitrarily long,
separating the adjectives and the noun they are modifying,
and has to involve at least some syntactic processes.

Another analysis is to utilize the fact that nouns in the construct chain 
are adjacent to each other:
probably, in Semitic languages there is a morphophonological rule 
prohibiting nouns with different definiteness from appearing together,
and if two nouns with the same definiteness appear together,
the first is realized as a construct state noun.
This is not a very local morphophonological rule,
but since Latin case agreement in noun phrases is also not very non-local this probably is acceptable. 

\bibliographystyle{plainnat}
\bibliography{arabic.bib}

\end{document}